\section{Results and interpretation}
\label{sec:results}
The principal equation is a fixed quartic in 238 unrestricted integer
variables, with 175 quadratic residuals, 5236 collected monomials and
coefficient height 3968. It removes six coordinates and one residual
from Round~19's balanced equation, at the same height and with 24 more
monomials. Against the newest Shared factorial error balanced equation,
it improves all four varying measures: variables, rows, support and height.
The 237-variable companion has a denser expansion; a separate 239-variable
factorial-free construction is sparser at the same height.

\begin{center}\small
\begin{tabular}{@{}lrrrr@{}}\toprule
Signed quartic & Variables & Rows & Terms & Height\\\midrule
Balanced principal &238&175&5236&3968\\
Minimum selected dimension &237&174&18569&7680\\
Power-base companion &239&175&5137&3968\\\bottomrule
\end{tabular}
\end{center}
All three have exact degree four over $\Z$. They are different tradeoffs,
not a single equation simultaneously optimal in every measure. The main
polynomial is fully defined by its printed residuals in
Appendix~\ref{sec:residuals} and its complete collected coefficient file.
It contains no factorial, variable exponent, binomial operation,
unspecified coefficient or external predicate.

The main arithmetic change replaces the shared bound
$M^2=1+4(S^2+\sum R_i^2)+G$ by
$M^2=2(S^2+\sum R_i^2)+G$.
Nonzeroness is recovered from the signed power request before division;
it is no longer a consequence of the norm alone. A radius-three-quarters
factorial lemma and the true remainder margin $1/32$ supply the required
uniqueness. The separate power-base lemma uses equality of factor counts
and positivity of all coefficients to recover bounded label products at
$F=t^t$. Neither argument simply assumes the signs or authenticity it
is meant to prove.

For height-one quartics, a fresh equality dummy $d$ can be replaced by
$a u$, where $u$ is an existing unit alias forced to one. The penalty
$C(a,b,au)$ is still nonnegative and has exactly the needed equality
zero set. Internal flow-edge contraction supplies another explicit
polynomial restoration. Every selected full polynomial is independently
collected to verify height one; the zero-set proof alone does not imply
that coefficient bound. The selected variable/support pairs are
\[
 \begin{gathered}
 (503,24193),\quad(505,20491),\quad(512,18255),\\
 (557,11999),\quad(617,8318),\quad(713,7176).
 \end{gathered}
\]
All five Round~19 selected points have a successor improving both dimension
and support, or dimension with unchanged support. The 503-variable point
is a separate smaller-dimension, larger-support tradeoff.

The new eleven-integer scalar points are exact degree/cost $298/147$
and $300/145$. The first improves both degree and cost over the last
$300/148$ delivery, but the newest incoming source already has degree
298 at 150 gates. The new achievement against that source is three fewer
gates. A second restricted degree theorem identifies 298 as the minimum
under constant radical weights and arbitrary integer polynomial centers
in that rebased frame; the old frame's distinct minimum remains 300.
Neither is a global lower bound on eleven-variable equations.

For each selected endpoint on its stated domain, the ordinary witness
proofs compose with the identified calibration to finite syntactic ZFC
refutation. Negating the existence equivalence gives unsolvability
provably equivalent in ZFC to $\Con(\mathrm{ZFC})$. This does not
prove consistency or determine whether any of the equations has a zero.
